\section{Auxiliary Graph-Learning Evaluation}
\label{sec:graph_evaluation}

\subsection{In-Distribution Ablations and Perturbation Asymmetry}
We first distinguish the class-balanced generator from the attribution and live-policy experiments. Their labels, model implementations, and observation units differ.

\paragraph{Skip-connection ablation}
The saved five-fold artifact reports Macro-F1 $0.9986 \pm 0.0029$ with skips and $0.4768 \pm 0.0239$ without skips. Recomputing the paired statistic from the five saved fold values gives $t=46.245$, $p=1.308\times10^{-6}$, and four degrees of freedom. The dispersions are population standard deviations across these folds. Shared training data and generator assumptions limit interpretation of this descriptive comparison. This ablation is not an empirical validation of the conditional derivative analysis in Appendix~\ref{sec:certgraph_arch}.

\paragraph{Class Proportions and Probability Interpretation}
The static experiments reported here are generator-based classification studies. Their balanced labels do not establish a 96\% enterprise safe base rate, and their softmax scores are not calibrated probabilities of real-world compromise. We retain probabilities and selected checkpoints in the matched audit so that classification errors can be checked independently.

\paragraph{Perturbation Asymmetry and Simple Comparators}
The saved robustness artifact gives Macro-F1 $0.9479$ after 30\% edge perturbation, compared with $0.8965$ and $0.7709$ after 5\% and 10\% feature perturbation. This sensitivity pattern motivates testing structural counterfactuals; it does not prove that every prediction ignores topology. A separately verified, previously saved three-seed ESC13 abstraction also favors simple graph-count models over the existing GNN on its IID cases. Its symbolic reference defines the labels, so its perfect agreement is a consistency result rather than independent exploit validation.

\subsection{Generalization and Hard-Negative Diagnostics}
A source audit identifies two problems with interpreting the original-label $41/63$ safe-label agreement as a general hard-negative detection result. Ten ``deep horizon'' cases retain a membership path to an allowed enrollment group and therefore are positive under a transitive effective-permission interpretation. The stored 22 disagreements are eight administrator-only write cases, four deny-write cases, seven benign OID links, and three deep OID links; the disagreements are reported using their stored categories. The canonical generator also repeats inputs across training and validation. These original-label counts are retained in the artifact but do not establish independent-case specificity or a binomial confidence interval.

We therefore freeze and execute a matched $2\times2$ diagnostic: additive versus feature-channel-gated implementations, each with counterfactual augmentation probability zero or 0.3, across three independently selected seeds. Model equations, propagation depth, and original artifacts remain unchanged. Each seed evaluates 250 opposite-label families, with membership depths 0, 1, 2, 4, and 8, randomized node order, and matched configuration features. Positive members retain reachable enrollment; negative members either restrict enrollment to the administrator or add a reachable explicit deny. Queue traversal and Boolean closure independently agree on all test labels under this stated ESC1 permission abstraction. CA trust and resolved permissions are assumed; Windows exploitability is not the label definition.

% Generated from independently verified saved checkpoint predictions.
\begin{table}[tbp]
\centering
\caption{Matched two-step diagnostic: 1,500 paired cases per arm (750 positive, 750 negative). Correct classification is a percentage; false positives and negatives are counts. Repeated training archetypes limit this synthetic diagnostic.}
\label{tab:hard_negatives}
\evidencetable
\begin{tabular}{llrrr}
\toprule
Implementation & Aug. & Correct & \makecell{False\\positive} & \makecell{False\\negative} \\
\midrule
Additive & 0 & 50.00\% & 750 & 0 \\
Additive & 0.3 & 71.67\% & 425 & 0 \\
\makecell[l]{Feature-channel\\gate} & 0 & 70.00\% & 0 & 450 \\
\makecell[l]{Feature-channel\\gate} & 0.3 & 70.00\% & 0 & 450 \\
\bottomrule
\end{tabular}
\end{table}


The gate gives zero false positives but misses all 450 positive cases at depths 2, 4, and 8 across the three seeds. Adding the same augmentation to the additive implementation improves its pooled accuracy from 50\% to 71.67\%, compared with the gate's 70\%. Thus the audit does not establish superiority of gating over an augmentation-matched comparator. First-step membership transforms are fixed in these existing NumPy implementations, and their feature-channel gate is $\sigma(W_{\mathrm{gate}}\mu_{\mathrm{reach}}+b)$, rather than a literal implementation of the general bilinear edge scorer in Section~\ref{sec:certgraph_arch}. Initialization also differs between implementations. These distinctions prevent attributing the entire contrast to a bilinear operation.

Only 22 unique training inputs occur per seed, with 21 also present in validation. Consequently, validation selection in this diagnostic is not an independent-distribution generalization test. All 6,000 saved test predictions are reproduced from saved checkpoints, and their labels and aggregates are checked independently. A symbolic second tier must also examine unresolved or out-of-horizon cases: verifying only neural-positive candidates cannot recover these observed neural false negatives.

\paragraph{Literal bilinear realization and family-disjoint validation}
We address the implementation mismatch with a separate literal realization of the bilinear edge and self-skip equations: two relational-attention layers, four heads, 32 hidden channels, type-specific input projection, and fully trainable encoder, attention, skip, and gate parameters. The binary readout evaluates the stated ESC1 resolved-permission predicate, not observed Windows compromise. Additive and gated arms start with identical parameter values and receive matched augmentation. Neither symbolic closure nor outcome labels enter the forward pass. Gradient, node-relabeling, and padding checks validate implementation properties; the feature-channel implementations are retained separately.

The family-disjoint generator constructs varied reachable and unreachable membership branches and paired allow/deny interventions. For each of three independently selected seeds, 300 training and 100 validation families are separated from 150 families in each of IID, long-path, and larger-background test cohorts. Typed five-round Weisfeiler--Lehman input hashes and family IDs are disjoint across splits within a seed. This conservatively excludes isomorphic whole inputs under the encoded roles and edge types, without claiming disjoint local motifs. Long-path positive enrollment chains have membership depth 2, 4, or 8. The larger-background cohort adds isolated nodes to different generated families; it is not a matched causal estimate of graph-size effects or enterprise-scale generalization.

% Generated from verified literal bilinear checkpoint predictions.
\begin{table}[tbp]
\centering
\caption{Literal two-layer attention: mean accuracy (\%) $\pm$ population SD over three seeds. Each test cohort contains 150 paired families per seed; initialization and augmentation are matched.}
\label{tab:literal_bilinear}
\evidencetable
\begin{tabular}{llrrr}
\toprule
Model & Aug. & IID & \makecell{Long\\paths} & \makecell{Larger\\background} \\
\midrule
\makecell[l]{Additive\\attention} & 0 & $86.00 \pm 1.78$ & $50.00 \pm 0.00$ & $87.00 \pm 0.82$ \\
\makecell[l]{Additive\\attention} & 0.3 & $86.00 \pm 1.78$ & $50.00 \pm 0.00$ & $87.00 \pm 0.82$ \\
\makecell[l]{Bilinear\\attention} & 0 & $86.00 \pm 1.78$ & $50.00 \pm 0.00$ & $87.00 \pm 0.82$ \\
\makecell[l]{Bilinear\\attention} & 0.3 & $86.00 \pm 1.78$ & $50.00 \pm 0.00$ & $87.00 \pm 0.82$ \\
\bottomrule
\end{tabular}
\end{table}


Independent NetworkX reachability checks all 5,100 saved case labels, and checkpoint inference reproduces all 10,800 held-out predictions. The literal bilinear and additive models have identical classifications in every matched arm: IID accuracy averages $86.00\%\pm1.79\%$, larger-background accuracy $87.00\%\pm0.82\%$, and long-path accuracy 50\%. In each arm, all 450 positive long-path cases are missed. Some IID negatives also require deny-membership reasoning beyond the local horizon; 126 pooled false positives remain. Probability scores differ, but their reported calibration errors concern these synthetic cohorts only. This closes the literal-implementation and repeated-whole-input gaps while failing to establish classification superiority for bilinear gating. Exact verification of unresolved cases remains necessary under the stated permission contract.

